\documentclass[12pt]{article}
\usepackage{amsmath,amssymb,amsfonts}
\usepackage[margin=1in]{geometry}

\begin{document}

\section*{Problem 5.27}

\textbf{Problem:} Show that if $E_\lambda(A)$ suffers a discontinuous jump at an eigenvalue $\lambda$ of $A$, then
\[
[E_\lambda(A) - E_{\lambda-}(A)]\Phi = \sum_{\nu=1}^{\mu_\lambda}c_\nu\Phi_\nu^{(\lambda)},
\]
where $\mu_\lambda$ is the multiplicity of $\lambda$ and $\{\Phi_\nu^{(\lambda)}\}$, $\nu = 1, 2, \ldots, \mu_\lambda$, is a set of orthonormal eigenvectors of $A$ belonging to $\lambda$, and $\lambda_1$ and $\lambda_2$ are points on either side of $\lambda$ in the intervals of continuity of $E(\lambda)$; that is, $E(\lambda)$ is constant in the intervals $[\lambda_1, \lambda)$ and $(\lambda, \lambda_2]$. $A$ is any self-adjoint operator.

\bigskip
\textbf{Solution:}

\subsection*{Setup}

Let $A$ be a self-adjoint operator with spectral resolution $E_\lambda$. The spectral theorem gives:
\[
A = \int_{-\infty}^{\infty}\lambda\,dE_\lambda.
\]

If $\lambda_0$ is an eigenvalue of $A$ with multiplicity $\mu$, then $E_\lambda$ has a jump discontinuity at $\lambda_0$. Define the jump:
\[
\Delta E_{\lambda_0} = E_{\lambda_0} - E_{\lambda_0^-} = \lim_{\epsilon\to 0^+}[E_{\lambda_0} - E_{\lambda_0-\epsilon}].
\]

\subsection*{Proof that $\Delta E_{\lambda_0}$ is the projection onto the eigenspace}

\textbf{Step 1:} $\Delta E_{\lambda_0}$ is a projection operator.

Since $E_\lambda$ is a resolution of the identity, for $\epsilon > 0$ small enough:
\[
\Delta E_{\lambda_0} = E_{\lambda_0} - E_{\lambda_0-\epsilon}
\]
is a projection (difference of two projections where the range of the smaller is contained in the range of the larger).

\textbf{Step 2:} The range of $\Delta E_{\lambda_0}$ is the eigenspace of $A$ at $\lambda_0$.

For any $\Phi$ in the Hilbert space, let $\Psi = \Delta E_{\lambda_0}\Phi$. Then:
\[
A\Psi = \int_{-\infty}^{\infty}\lambda\,dE_\lambda(\Delta E_{\lambda_0}\Phi).
\]

Since $\Delta E_{\lambda_0}$ projects onto the subspace corresponding to the single point $\lambda_0$:
\[
dE_\lambda(\Delta E_{\lambda_0}\Phi) = \begin{cases} \Psi & \text{at } \lambda = \lambda_0, \\ 0 & \text{otherwise}. \end{cases}
\]

Therefore $A\Psi = \lambda_0\Psi$, showing that $\Psi$ is an eigenvector with eigenvalue $\lambda_0$.

\textbf{Step 3:} Conversely, if $A\Phi = \lambda_0\Phi$, then $\Delta E_{\lambda_0}\Phi = \Phi$.

From the spectral theorem:
\[
\|A\Phi - \lambda_0\Phi\|^2 = \int_{-\infty}^{\infty}(\lambda - \lambda_0)^2\,d\|E_\lambda\Phi\|^2 = 0.
\]

Since the integrand is non-negative, $d\|E_\lambda\Phi\|^2 = 0$ for $\lambda \neq \lambda_0$. This means all the ``spectral weight'' of $\Phi$ is concentrated at $\lambda_0$, so $\Delta E_{\lambda_0}\Phi = \Phi$.

\textbf{Step 4:} Expanding in the eigenbasis.

Since $\Delta E_{\lambda_0}$ projects onto the $\mu_\lambda$-dimensional eigenspace spanned by the orthonormal eigenvectors $\{\Phi_\nu^{(\lambda)}\}_{\nu=1}^{\mu_\lambda}$:
\[
\Delta E_{\lambda_0}\Phi = \sum_{\nu=1}^{\mu_\lambda}(\Phi_\nu^{(\lambda)}, \Phi)\,\Phi_\nu^{(\lambda)} = \sum_{\nu=1}^{\mu_\lambda}c_\nu\Phi_\nu^{(\lambda)},
\]
where $c_\nu = (\Phi_\nu^{(\lambda)}, \Phi)$.

\[
\boxed{[E_{\lambda_0} - E_{\lambda_0^-}]\Phi = \sum_{\nu=1}^{\mu_\lambda}c_\nu\Phi_\nu^{(\lambda)}, \quad c_\nu = (\Phi_\nu^{(\lambda)}, \Phi).}
\]

\end{document}
